\section{Вариационные принципы построения математических моделей}
\label{sec:variational-principles-of-modeling}

\textbf{Вариационный принцип} --- способ построения математической модели,
при котором искомое состояние или траектория выделяются среди допустимых
состояний условием стационарности некоторого функционала. При этом
необходимо заранее указать класс допустимых функций и ограничения задачи.

\textbf{Функционал} ставит функции в соответствие число. Типичный пример:
\[
J[u]=\int_a^b F(x,u,u')\,dx.
\]
Условие стационарности записывается как
\[
\boxed{\delta J[u]=0}.
\]
В общем случае стационарная точка не обязана быть минимумом: она может быть
максимумом или седловой точкой.

\subsection{Первая вариация и уравнение Эйлера--Лагранжа}

Пусть $u\in C^2[a,b]$, а вариации $\eta$ удовлетворяют
$\eta(a)=\eta(b)=0$. Рассмотрим семейство
\[
u_\varepsilon=u+\varepsilon\eta.
\]
Первая вариация определяется как
\[
\delta J[u;\eta]
=\left.\frac{d}{d\varepsilon}J[u+\varepsilon\eta]\right|_{\varepsilon=0}.
\]
Для функционала $J[u]=\int_a^b F(x,u,u')\,dx$ после интегрирования по частям
получаем
\[
\delta J
=\int_a^b
\left(
\frac{\partial F}{\partial u}
-\frac{d}{dx}\frac{\partial F}{\partial u'}
\right)\eta\,dx.
\]

Центральным результатом является \textbf{основная лемма вариационного
исчисления}: если непрерывная функция $g$ удовлетворяет
\[
\int_a^b g(x)\eta(x)\,dx=0
\]
для всех гладких функций $\eta$ с компактным носителем в $(a,b)$, то
$g(x)\equiv0$. Поэтому необходимое условие стационарности имеет вид
\[
\boxed{
\frac{\partial F}{\partial u}
-\frac{d}{dx}\left(\frac{\partial F}{\partial u'}\right)=0
}
\]
--- это \textbf{уравнение Эйлера--Лагранжа}.

\subsection{Принцип стационарного действия}

Для консервативной механической системы действие равно
\[
\boxed{S[q]=\int_{t_1}^{t_2}L(q,\dot q,t)\,dt},
\qquad L=T-V,
\]
где $T$ и $V$ --- кинетическая и потенциальная энергии. Принцип Гамильтона
утверждает, что действительное движение при фиксированных концах
\[
\delta q(t_1)=\delta q(t_2)=0
\]
делает действие стационарным:
\[
\boxed{\delta S=0}.
\]
Отсюда для обобщённых координат $q_i$ следуют уравнения Лагранжа
\[
\boxed{
\frac{d}{dt}\frac{\partial L}{\partial\dot q_i}
-\frac{\partial L}{\partial q_i}=0,
\qquad i=1,\ldots,n.
}
\]
Таким образом, один интегральный принцип порождает систему
дифференциальных уравнений движения.

\subsection{Принцип возможных перемещений и минимум потенциальной энергии}

\textbf{Виртуальным перемещением} называют бесконечно малое изменение
конфигурации, совместимое со связями в данный момент времени. Для системы
с идеальными связями реакции связей не совершают виртуальной работы.
Условие равновесия можно записать как
\[
\boxed{\sum_i \mathbf F_i\cdot\delta\mathbf r_i=0}
\]
для всех допустимых виртуальных перемещений, где $\mathbf F_i$ --- активные
силы.

Для консервативной упругой системы вводится полная потенциальная энергия
\[
\boxed{\Pi=U-W},
\]
где $U$ --- энергия деформации, $W$ --- потенциал работы внешних сил в
принятой знаковой конвенции. Равновесие удовлетворяет
\[
\delta\Pi=0.
\]
Если вторая вариация положительна на ненулевых допустимых вариациях, то
равновесие устойчиво и реализует локальный минимум $\Pi$.

\subsection{Связи и множители Лагранжа}

Пусть на движение наложены голономные связи
\[
f_\alpha(q,t)=0,\qquad \alpha=1,\ldots,m.
\]
Их удобно учитывать с помощью множителей Лагранжа $\lambda_\alpha(t)$,
заменяя лагранжиан расширенным:
\[
L^*(q,\dot q,\lambda,t)
=L(q,\dot q,t)
+\sum_{\alpha=1}^m\lambda_\alpha(t)f_\alpha(q,t).
\]
Стационарность действия с $L^*$ даёт одновременно уравнения движения и
уравнения связей. Множители часто имеют физический смысл реакций связей.
Аналогичный приём применяется и в статических вариационных задачах с
ограничениями.

\subsection{Диссипативные системы}

Для линейной вязкой диссипации вводят \textbf{диссипативную функцию Рэлея}
\[
\mathcal R=\frac12\sum_{i,j}c_{ij}\dot q_i\dot q_j.
\]
Если $Q_i$ обозначают остальные непотенциальные обобщённые силы, то
уравнения Лагранжа принимают вид
\[
\boxed{
\frac{d}{dt}\frac{\partial L}{\partial\dot q_i}
-\frac{\partial L}{\partial q_i}
+\frac{\partial\mathcal R}{\partial\dot q_i}
=Q_i.
}
\]
Например, для осциллятора с вязким сопротивлением
\[
L=\frac{m\dot x^2}{2}-\frac{kx^2}{2},
\qquad \mathcal R=\frac{c\dot x^2}{2},
\]
получается
\[
m\ddot x+c\dot x+kx=F(t).
\]

\subsection{Пример: упругая балка}

Для балки Эйлера--Бернулли с поперечным перемещением $u(x)$ функционал
полной потенциальной энергии можно записать как
\[
\Pi[u]
=\frac12\int_0^L EI\,(u'')^2\,dx
-\int_0^L q(x)u(x)\,dx.
\]
Условие $\delta\Pi=0$ после двукратного интегрирования по частям приводит к
\[
EIu^{(4)}=q(x)
\]
и соответствующим естественным и существенным граничным условиям.
Именно эта связь вариационной и дифференциальной постановок лежит в
основе метода Ритца и метода конечных элементов.


\subsection{Вторая вариация и различие стационарности и минимума}

Уравнение Эйлера--Лагранжа даёт только необходимое условие стационарности.
Чтобы различить минимум, максимум и седловую точку, исследуют вторую вариацию.
Для
\[
    J[u]=\int_a^b F(x,u,u')\,dx
\]
при фиксированных концах
\[
\delta^2J[u;\eta]
=\int_a^b\left(
F_{uu}\eta^2+2F_{u u'}\eta\eta'+F_{u'u'}(\eta')^2
\right)dx,
\]
где производные $F$ вычисляются на стационарной функции $u$. Для локального
минимума необходимо $\delta^2J[u;\eta]\ge0$ для всех допустимых $\eta$.
Одним из классических необходимых условий является условие Лежандра
\[
    F_{u'u'}\ge0,
\]
а усиленное условие $F_{u'u'}>0$ входит в стандартные достаточные условия
сильного локального минимума вместе с дополнительными условиями Якоби и
Вейерштрасса. Поэтому при построении модели важно не отождествлять
\(\delta J=0\) и принцип минимума без дополнительных предпосылок.

\subsection{Граничные члены и естественные условия}

При интегрировании по частям первая вариация имеет более полный вид
\[
\delta J
=\int_a^b\left(F_u-\frac{d}{dx}F_{u'}\right)\eta\,dx
+\left[F_{u'}\eta\right]_a^b.
\]
Если значение $u$ на конце закреплено, соответствующая вариация равна нулю.
Если же конец свободен, из произвольности вариации возникает
\textbf{естественное граничное условие}. Например, при свободном значении
$u(b)$ получаем
\[
    F_{u'}(b)=0.
\]
В задачах механики сплошной среды та же логика приводит к разделению
\textbf{существенных} граничных условий, которые включаются в класс допустимых
функций, и \textbf{естественных} условий, которые возникают из граничного члена
вариационной формулы.

\subsection{Вариационная постановка краевой задачи: принцип Дирихле}

Рассмотрим эллиптическую задачу
\[
    -\nabla\cdot(k\nabla u)=f\quad\text{в }\Omega,
    \qquad u=0\quad\text{на }\Gamma_D,
    \qquad k\frac{\partial u}{\partial n}=g\quad\text{на }\Gamma_N,
\]
где $k(x)\ge k_0>0$. Ей соответствует функционал
\[
    J[v]
    =\frac12\int_\Omega k|\nabla v|^2\,dx
     -\int_\Omega fv\,dx
     -\int_{\Gamma_N}gv\,ds
\]
на пространстве
\[
    V=\{v\in H^1(\Omega):v|_{\Gamma_D}=0\}.
\]
Условие $\delta J[u;w]=0$ для всех $w\in V$ даёт слабую форму
\[
    \boxed{
    \int_\Omega k\nabla u\cdot\nabla w\,dx
    =\int_\Omega fw\,dx+\int_{\Gamma_N}gw\,ds
    \quad\forall w\in V.
    }
\]
Если $\Gamma_D$ имеет ненулевую меру, то неравенство Пуанкаре вместе с
$k\ge k_0$ обеспечивает коэрцитивность. Поэтому функционал строго выпуклый,
и стационарная точка одновременно является единственным минимумом. Это пример,
когда вариационный принцип не только компактно задаёт модель, но и сразу даёт
естественную функционально-аналитическую постановку.

\subsection{Методы Ритца и Галёркина как следствие вариационного принципа}

Пусть $V_N=\operatorname{span}\{\varphi_1,\ldots,\varphi_N\}\subset V$ и
\[
    u_N=\sum_{j=1}^N a_j\varphi_j.
\]
В методе Ритца минимизируют $J[u_N]$ по коэффициентам $a_j$:
\[
    \frac{\partial J}{\partial a_i}=0,
    \qquad i=1,\ldots,N.
\]
Для квадратичного функционала это приводит к системе
\[
    \sum_{j=1}^N a(u_N,\varphi_i)=\ell(\varphi_i),
\]
то есть к галёркинской ортогональности невязки. В симметричной коэрцитивной
задаче методы Ритца и Галёркина совпадают. Именно поэтому вариационные принципы
являются естественным мостом от физических законов к вычислительным моделям и
методу конечных элементов.

\subsection{Что важно сказать на экзамене}

Нужно сформулировать понятия функционала и первой вариации, основную
лемму вариационного исчисления и уравнение Эйлера--Лагранжа. В механике
ключевые примеры --- принцип стационарного действия, принцип возможных
перемещений и стационарность полной потенциальной энергии. Следует
помнить, что стационарность не всегда означает минимум, а связи можно
включать с помощью множителей Лагранжа.

\newpage
